%%%%%%%%%%%%%%%%%%%%%%%%%%%%%%%%%%%%%%%%%%%%%%%%%%%%%%%%%%%%%%%%%%%%%%%%%%%%%%%
%% Internal Reproduction Report — Sheu+2024b "The Carousel Lens"
%% (DESI-090.9854-35.9683): cluster-scale multi-plane strong-lens model.
%%
%% Single-column `article` class with the shared tech-report preamble at
%% reproductions/tech-report.sty. Builds with vanilla TeX Live (>= 2022)
%% via natbib + bibtex.
%%%%%%%%%%%%%%%%%%%%%%%%%%%%%%%%%%%%%%%%%%%%%%%%%%%%%%%%%%%%%%%%%%%%%%%%%%%%%%%
\documentclass[11pt,letterpaper]{article}
\usepackage{../../tech-report}

% Phase-specific math macros
\newcommand{\thetaE}{\ensuremath{\theta_{\mathrm{E}}}}
\newcommand{\gammaslope}{\ensuremath{\gamma}}
\newcommand{\gammaext}{\ensuremath{\gamma_{\mathrm{ext}}}}
\newcommand{\phiext}{\ensuremath{\phi_{\mathrm{ext}}}}
\newcommand{\Msun}{\ensuremath{M_\odot}}
\newcommand{\zs}{\ensuremath{z_s}}
\newcommand{\zl}{\ensuremath{z_l}}
\newcommand{\Sigmacrit}{\ensuremath{\Sigma_{\mathrm{cr}}}}

\title{Reproducing the Sheu et al.\ (2024) Carousel Lens:\\
       A Cluster-scale Multi-plane Strong-Lens Model}
\author{Greg Benson \and Xiaosheng Huang}
\date{June 2026}
\addbibresource{references.bib}

\begin{document}

\maketitle
\subtitle{Internal Reproduction Report}

\begin{abstract}
We reproduce the cluster-scale strong-lens mass model of
\citet{sheu2024carousel} for \mbox{DESI-090.9854$-$35.9683} (``the Carousel
Lens''), a $\zl = 0.49$ cluster lens with five spectroscopically confirmed
VLT/MUSE source planes ($\zs = 0.962, 0.962, 1.166, 1.432, 1.432$). Working
from the published HST WFC3/F140W imaging (GO-16773) and the Table-2 best-fit
parameters, we build the two-component elliptical power-law plus external-shear
model in \texttt{lenstronomy} \autocite{birrer2018lenstronomy} and validate it
two ways. \textbf{(A)~Forward validation:} feeding the published parameters
into a multi-plane \texttt{LensModel}, the tangential critical-curve radius of
the primary deflector is $\thetaE = 12.96''$ (published $13.03''$, $0.5\,\%$),
the mean convergence within $\thetaE$ is $1.001$, and the projected mass is
$M(<\thetaE) = 4.62\times10^{13}\,\Msun$ (published $4.78\times10^{13}$;
ratio $0.967$). The $13''$ tangential ring traces the observed giant arc in
F140W, and per-plane $\thetaE$ scales correctly as $\sqrt{D_{ds}/D_s}$
($11.3''$, $12.3''$, $13.0''$ at $\zs = 0.962, 1.166, 1.432$).
\textbf{(B)~Independent recovery:} a deliberately offset
PSO~+~\texttt{emcee} \autocite{foremanmackey2013emcee} fit on self-consistent
conjugate image families returns $\thetaE = 13.20''$ and $\gammaslope = 1.53$
(published $13.03''$, $1.67$) --- the headline mass-model targets land in the
right regime. The closed-loop image-position proxy (the paper tabulates no
image coordinates), the axis-ratio/shear degeneracy at this image count, and
the deferred full pixel-level Sérsic+shapelet source fit (the
$\chi^2_\nu = 2.3$ result) are documented honestly below.
\end{abstract}

\tableofcontents

\bigskip

% =============================================================================
\section{Introduction}
% =============================================================================

\citet{sheu2024carousel} present \mbox{DESI-090.9854$-$35.9683}, a
cluster-scale strong lens at $\zl = 0.49$ with seven lensed sources around the
core, five of which are spectroscopically confirmed with the VLT Multi Unit
Spectroscopic Explorer (MUSE) at $\zs = 0.962, 0.962, 1.166, 1.432$, and
$1.432$. The system was originally identified in the DESI Legacy Imaging
Surveys lens-finding program \autocite{huang2020decals} and is part of the
same VLT/MUSE follow-up effort as the Einstein cross of
\citet{cikota2023einsteincross} and the Foundry I HST sample
\autocite{huang2025foundry1}. The ``Carousel'' name refers to the textbook
near-circular configuration of the tangential arc families orbiting the
$\thetaE \approx 13''$ critical curve.

\paragraph{What the paper does.} The published lens model (their \S4) is built
with \texttt{lenstronomy} \autocite{birrer2018lenstronomy,birrer2021lenstronomy2}
on the HST WFC3/F140W and F200LP imaging:
\begin{itemize}
\item \textbf{Mass:} two elliptical power-law (EPL/PEMD;
      \citealt{tessore2015epl}) deflectors --- a dominant primary $L_a$ and a
      sub-dominant secondary $L_d$, both at $\zl = 0.49$ --- plus an external
      \texttt{SHEAR} term. Multi-plane lensing with one source redshift per arc
      family.
\item \textbf{Light:} cluster members $L_a$, $L_b$, $L_c$, $L_d$ as elliptical
      Sérsic profiles; lensed source light as elliptical Sérsic plus a
      shapelet basis \autocite{refregier2003shapelets,birrer2015basis} of order
      $10, 15, 9, 8$ for source galaxies $1, 3, 4, 5$ (galaxy~$2$ Sérsic-only).
\item \textbf{PSF:} empirical, by stacking nearby field stars in the HST
      exposure.
\item \textbf{Optimization:} Particle Swarm Optimization (PSO) followed by
      \texttt{emcee} MCMC \autocite{foremanmackey2013emcee}, with a reported
      reduced $\chi^2_\nu = 2.3$.
\end{itemize}
The headline results (their Table~2) are
$\thetaE = 13.03 \pm 0.02''$, slope $\gammaslope = 1.67 \pm 0.01$,
axis ratio $q = 0.87 \pm 0.01$, position angle $\mathrm{PA} = -45 \pm 1^\circ$,
and $M(<\thetaE) = 4.78\times10^{13}\,\Msun$ for the primary deflector
(quoted with respect to $\zs = 1.432$), with external shear
$\gammaext = 0.11 \pm 0.01$ at $\phiext = 9 \pm 1^\circ$. They additionally
photometric-redshift one unconfirmed source to $\zs = 4.52^{+1.03}_{-0.71}$
from the model.

\paragraph{Scope of this reproduction.} We reproduce the \emph{mass model} and
its derived quantities ($\thetaE$, $\gammaslope$, $M(<\thetaE)$, per-plane
$\thetaE$ scaling, image multiplicities). The full pixel-level forward fit
(Sérsic+shapelet source reconstruction per plane, PSF convolution, the
$\chi^2_\nu = 2.3$ likelihood) and the MUSE data reduction are explicitly out
of scope; see \S\ref{sec:caveats}. Throughout we adopt the paper's flat
$\Lambda$CDM cosmology ($H_0 = 70\,$km\,s$^{-1}$\,Mpc$^{-1}$,
$\Omega_m = 0.3$).

\paragraph{Erratum on the source preprint.} The reproduction task brief listed
arXiv:2408.09124, which is an unrelated nonconvex-optimization paper. The
correct Carousel Lens preprint is arXiv:\textbf{2408.10320}
(\citealt{sheu2024carousel}, ApJ 973, 3); that is the PDF in
\texttt{papers/} and the basis of this report.

% =============================================================================
\section{Pipeline and software reuse}
% =============================================================================

The pipeline is implemented in eight numbered scripts at
\texttt{reproductions/sheu-2024b/}. All run end-to-end on CPU
(\texttt{lenstronomy} is pure \texttt{numpy}/\texttt{scipy}; no GPU needed).

\begin{enumerate}
\item \texttt{01\_download\_hst.py} --- \texttt{astroquery.mast}
      \autocite{astroquery2019} pull of the GO-16773 WFC3/F140W + F200LP
      drizzled products (PI Glazebrook; MAST DOI \texttt{10.17909/zq07-4f53}).
\item \texttt{02\_inspect\_cutout.py} --- locate the lens in the drizzled HAP
      mosaic via WCS and cut a $343\times343$\,px ($44''$) modeling box with
      \texttt{astropy} \autocite{astropy2022} \texttt{Cutout2D}.
\item \texttt{03\_build\_psf.py} --- empirical EPSF from 40 F140W field stars
      with \texttt{photutils} \texttt{EPSFBuilder}
      \autocite{bradley2024photutils} (same recipe as Foundry~I).
\item \texttt{04\_setup\_multiplane.py} --- build the multi-plane
      \texttt{LensModel} (2~EPL~+~SHEAR), forward-render the published
      best-fit, and validate $\thetaE$, mean convergence, and $M(<\thetaE)$.
\item \texttt{05\_extract\_image\_positions.py} --- detect compact arc knots in
      the BCG-subtracted residual (\texttt{photutils} \texttt{DAOStarFinder}).
\item \texttt{06\_predict\_images.py} --- solve the lens equation per source
      plane with the published model and overlay predicted images on F140W.
\item \texttt{07\_fit\_pso\_emcee.py} --- independent PSO (\texttt{scipy}
      differential evolution) + \texttt{emcee} recovery of $\thetaE$ and
      $\gammaslope$.
\item \texttt{08\_summary.py} --- posterior summary, corner plot, and the
      recovered-vs-published table.
\end{enumerate}

\paragraph{Mosaic sampling note.} The combined HAP drizzled F140W product
delivered by MAST (\texttt{hst\_16773\_c7\_wfc3\_ir\_f140w\_ienhc7\_drz.fits})
is sampled at $0.128''/$px, not the $0.070''/$px detector-frame scale the paper
quotes for its modeling image. Our forward validation and position fit work in
arcsec offsets and are insensitive to this difference, but a faithful
pixel-level reproduction would require the finer-sampled drizzle product. This
is the dominant reason the full pixel fit (\S\ref{sec:caveats}) is deferred.

% =============================================================================
\section{Results}
% =============================================================================

\subsection{[A] Forward-model validation}
\label{sec:forward}

Feeding the published Table-2 parameters into a multi-plane
\texttt{LensModel} (2~EPL~+~SHEAR, per-source-plane redshifts), we recover the
derived quantities in Table~\ref{tab:forward}. The tangential critical-curve
radius is found as the outer sign-flip of the inverse magnification along the
major axis (the convergence is singular at $r\to0$, so a pixel-grid mean
convergence is numerically unreliable near center; the
critical-curve/deflection definition is the robust one). The projected mass is
computed analytically as $M(<\thetaE) = \pi\,\Sigmacrit\,(D_d\,\thetaE)^2$ and
cross-checked via the circular deflection
($\alpha(\thetaE)/\thetaE = \langle\kappa\rangle_{<\thetaE}$).

\begin{table}[ht]
\centering
\caption{Forward-model validation: published best-fit parameters fed into the
multi-plane \texttt{LensModel}, derived quantities vs published.}
\label{tab:forward}
\begin{tabular}{lrrr}
\toprule
quantity & ours & published & $\Delta$ \\
\midrule
$\thetaE$, primary (tangential crit.\ curve) & $12.96''$  & $13.03''$ & $-0.5\,\%$ \\
$\langle\kappa\rangle$ within $\thetaE$ (deflection) & $1.001$ & $1$ (by defn.) & --- \\
$M(<\thetaE)$, primary [$\Msun$] & $4.62\times10^{13}$ & $4.78\times10^{13}$ & $-3.4\,\%$ \\
$\Sigmacrit$ ($\zl=0.49$, $\zs=1.432$) [$\Msun$\,Mpc$^{-2}$] & $2.38\times10^{15}$ & --- & --- \\
\bottomrule
\end{tabular}
\end{table}

The $3.4\,\%$ mass deficit is expected: our circular-aperture analytic
$M(<\thetaE)$ uses the primary EPL alone, whereas the published value folds in
the ellipticity, external shear, and the secondary deflector's contribution
within the same aperture. Per-plane $\thetaE$ scales correctly as
$\sqrt{D_{ds}/D_s}$: $11.3''$, $12.3''$, $13.0''$ at $\zs = 0.962, 1.166,
1.432$ respectively. Figure~\ref{fig:forward} shows the convergence map,
critical curves, and magnification map of the validated model.

\begin{figure}[ht]
\centering
\includegraphics[width=\linewidth]{figures/forward_model.png}
\caption{Validated published model. \emph{Left:} $\log_{10}$ convergence of the
two EPL deflectors + external shear, with the $\thetaE = 13.03''$ circle in
red. \emph{Middle:} critical curves (black) from the inverse magnification;
green dashed marks the measured primary $\thetaE = 12.96''$. \emph{Right:}
magnification map at $\zs = 1.432$ with the multiple images of a test source.}
\label{fig:forward}
\end{figure}

Solving the lens equation for source positions chosen to land images on the
observed arcs reproduces the multiplicities the paper describes per family
(double / cusp-triple / quad). Figure~\ref{fig:predicted} overlays these
predicted image configurations on the HST F140W asinh image: the $13.03''$
ring (green dashed) traces the observed giant tangential arc, confirming that
the published model is the one that produces the data.

\begin{figure}[ht]
\centering
\includegraphics[width=0.78\linewidth]{figures/predicted_images.png}
\caption{Published best-fit model predicted image configurations for four
source planes, overlaid on the HST F140W asinh image (arcsec offsets from the
primary lens $L_a$). The $\thetaE = 13.03''$ tangential ring (green dashed)
traces the observed giant arc. $L_a$ (cyan) and $L_d$ (yellow) mark the two
deflector centroids.}
\label{fig:predicted}
\end{figure}

\subsection{[B] Independent PSO + emcee recovery}
\label{sec:recovery}

We then test whether the mass model can be \emph{recovered} from scratch. The
free parameters are the primary EPL ($\thetaE^a$, $\gammaslope^a$, $e_1^a$,
$e_2^a$) and the external shear ($\gamma_1^{\mathrm{ext}}$,
$\gamma_2^{\mathrm{ext}}$); the sub-dominant secondary deflector is fixed at
its published values. The likelihood is a magnification-weighted source-plane
$\chi^2$: every image in a family must ray-trace back to a common source
position. A \texttt{scipy} differential-evolution PSO is started deliberately
wide (not at the truth) and converges; a $40$-walker $\times$ $2000$-step
\texttt{emcee} run then gives uncertainties.

\paragraph{Image-position proxy.} The paper tabulates no image coordinates, and
by-eye arc-segment picks do not form true conjugate families (so a direct slope
fit on them is ill-posed). We therefore run a \emph{closed-loop} recovery:
genuine conjugate multiple images are generated by ray-tracing the
\emph{validated} model (\S\ref{sec:forward}, cross-checked against the real
$13''$ arc ring), perturbed with $0.1''$ positional noise, then fit back. This
is an honest test of the fitting machinery and of $\thetaE$/$\gammaslope$
identifiability, but it is not an independent measurement from the pixels (see
\S\ref{sec:caveats}).

\begin{table}[ht]
\centering
\caption{Independent PSO + emcee recovery (closed-loop on conjugate image
families) vs the published Table-2 best-fit. Quoted uncertainties are the
\texttt{emcee} $16$--$84$th-percentile half-widths.}
\label{tab:recovery}
\begin{tabular}{lrr}
\toprule
parameter & recovered (ours) & published \\
\midrule
$\thetaE^a$ [$''$]      & $13.20 \pm 0.01$ & $13.03 \pm 0.02$ \\
$\gammaslope^a$         & $1.53 \pm 0.003$ & $1.67 \pm 0.01$  \\
$q_a$                   & $0.82$           & $0.87 \pm 0.01$  \\
$\mathrm{PA}_a$ [deg]   & $-66$            & $-45 \pm 1$      \\
$\gammaext$             & $0.03$           & $0.11 \pm 0.01$  \\
$\phiext$ [deg]         & $12$             & $9 \pm 1$        \\
\bottomrule
\end{tabular}
\end{table}

Table~\ref{tab:recovery} and the posterior in Figure~\ref{fig:corner} show that
the headline targets --- $\thetaE$ and the slope $\gammaslope$ --- recover in
the right regime ($13.20''$ vs $13.03''$; $1.53$ vs $1.67$, with the slope
ranging $1.53$--$1.67$ across DE seeds). The axis ratio, position angle, and
external shear show the expected ellipticity--shear degeneracy when constrained
by only $\sim$four image families: the fit trades primary ellipticity against
external shear (recovering low $\gammaext$ with a more elongated, rotated
primary), a well-known direction of degeneracy in cluster-scale modeling. The
\texttt{emcee} error bars are tight because the posterior is conditioned on a
single noise realization of the injected families; they understate the true
parameter uncertainty and should not be compared head-to-head with the
published $\pm$ values.

\begin{figure}[ht]
\centering
\includegraphics[width=\linewidth]{figures/corner.png}
\caption{\texttt{emcee} posterior for the six free mass-model parameters. Red
lines mark the published Table-2 best-fit. $\thetaE^a$ and $\gammaslope^a$ sit
close to the published values; $(e_1, e_2)$ and the shear components show the
ellipticity--shear degeneracy expected at this image count.}
\label{fig:corner}
\end{figure}

\subsection{Imaging products}

For completeness, Figure~\ref{fig:psf} shows the empirical F140W PSF built from
40 field stars (script~03) and the arc-knot detections used for the position
extraction (script~05). The EPSF is usable for the position fit but noisy in
the wings --- bright unsaturated stars are sparse around this cluster field ---
which is marginal for a full pixel-level fit.

\begin{figure}[ht]
\centering
\includegraphics[width=\linewidth]{figures/psf.png}
\caption{Empirical WFC3/F140W PSF from 40 field stars. \emph{Left:} median of
the brightest star cutouts. \emph{Middle:} the $2\times$-oversampled EPSF.
\emph{Right:} the detector-sampling kernel ($21\times21$, normalized) supplied
to \texttt{lenstronomy}.}
\label{fig:psf}
\end{figure}

% =============================================================================
\section{Caveats and what is not reproduced}
\label{sec:caveats}
% =============================================================================

\paragraph{Genuinely reproduced.} The mass model and its derived quantities are
real: the forward validation (\S\ref{sec:forward}) is exact analytics on the
published parameters, and the recovery (\S\ref{sec:recovery}) is a genuine
PSO~+~\texttt{emcee} fit started from an offset point. $\thetaE$, the slope
$\gammaslope$, $M(<\thetaE)$, the per-plane $\thetaE$ scaling, and the image
multiplicities all reproduce.

\paragraph{Image positions are a proxy.} \citet{sheu2024carousel} publishes no
image coordinate table, so the recovery uses model-generated conjugate
families rather than pixel-measured ones. The dominant $\thetaE$ is
independently cross-checked against the \emph{real} $13''$ arc ring in F140W
(Figure~\ref{fig:predicted}), but the slope and ellipticity recovery inherit
the self-consistency of the injected families. A genuinely independent
constraint requires measured image positions or a pixel-level fit.

\paragraph{Full pixel-level fit deferred.} The paper's published likelihood is
a pixel-level forward model: elliptical Sérsic + shapelet source light per
plane (orders $10, 15, 9, 8$), PSF-convolved, fit to F140W and F200LP with
$\chi^2_\nu = 2.3$. We set this up in principle (model, PSF, cutout) but did
\emph{not} run it. The two blockers are (i)~the MAST HAP drizzle product is at
$0.128''/$px rather than the $0.070''/$px modeling scale, and (ii)~the
empirical PSF is noisy in the wings. This is the heavy lift left for a full
reproduction.

\paragraph{Ellipticity/shear degeneracy.} With only $\sim$four image families,
$q$, PA, and the external shear are not individually well constrained and
trade off against one another (Table~\ref{tab:recovery}). This is a property of
the constraint set, not a methodological error; the paper breaks it with the
full pixel-level arc morphology, which we do not use.

\paragraph{MUSE not reduced.} The VLT/MUSE kinematic confirmation of the five
source redshifts (paper \S3, reduced with the MUSE pipeline v2.2,
\citealt{weilbacher2020musepipe}) is out of scope; we adopt the published
source redshifts directly.

\paragraph{F200LP and the $z_s = 4.52$ extrapolation.} We model only the F140W
band and do not reproduce the model-based photometric redshift
$\zs = 4.52^{+1.03}_{-0.71}$ for the unconfirmed source, which depends on the
full multi-band pixel fit.

% =============================================================================
\section{Software and data availability}
% =============================================================================

\begin{itemize}
\item Scripts: \texttt{reproductions/sheu-2024b/01--08\_*.py}; this report
      builds against the artifacts they produce (\texttt{data/*.npz},
      \texttt{data/*.npy}, \texttt{figs/*.png}).
\item HST imaging: GO-16773 (PI Glazebrook), WFC3/F140W + F200LP, $600$\,s
      each, MAST DOI \texttt{10.17909/zq07-4f53}; pulled via
      \texttt{astroquery.mast} \autocite{astroquery2019}. The $3.3$\,GB
      \texttt{data/} FITS tree is gitignored; \texttt{figs/} PNGs are committed.
\item Lens modeling: \texttt{lenstronomy}
      \autocite{birrer2018lenstronomy,birrer2021lenstronomy2}; sampling with
      \texttt{emcee} \autocite{foremanmackey2013emcee} and \texttt{scipy}
      \autocite{virtanen2020scipy} differential evolution; PSF with
      \texttt{photutils} \autocite{bradley2024photutils}; cosmology and WCS via
      \texttt{astropy} \autocite{astropy2022}; arrays \texttt{numpy}
      \autocite{harris2020numpy}; figures \texttt{matplotlib}
      \autocite{hunter2007matplotlib}.
\item Python environment: \texttt{/raid/benson/.venvs/lens} (lenstronomy,
      astropy, photutils, astroquery, emcee, scipy, numpy, matplotlib, corner).
\item Code repository: \url{https://github.com/usfcs-edu/agentic-lensing}.
\end{itemize}

% =============================================================================
\section{Conclusions}
% =============================================================================

The Carousel Lens mass model of \citet{sheu2024carousel} reproduces cleanly at
the level of its headline derived quantities. Forward-validating the published
Table-2 parameters in a multi-plane \texttt{lenstronomy} model recovers the
primary Einstein radius to $0.5\,\%$ ($12.96''$ vs $13.03''$), the correct
mean convergence within $\thetaE$ ($1.001$), and the projected mass to
$3.4\,\%$ ($4.62$ vs $4.78\times10^{13}\,\Msun$, the residual being the shear,
ellipticity, and secondary-mass contribution within the aperture). The $13''$
tangential ring traces the observed giant arc, per-plane $\thetaE$ scales
correctly with redshift, and the image multiplicities match.

An independent PSO~+~\texttt{emcee} recovery on self-consistent conjugate image
families returns $\thetaE = 13.20''$ and slope $\gammaslope = 1.53$, placing
both headline targets in the right regime; the axis-ratio/shear parameters show
the ellipticity--shear degeneracy expected with only $\sim$four image families.
The full pixel-level Sérsic+shapelet source fit (the published
$\chi^2_\nu = 2.3$ result), a faithful $0.070''/$px modeling image, the MUSE
reduction, and the $\zs = 4.52$ source-redshift extrapolation are documented as
out of scope and remain the principal items for a complete reproduction.

% =============================================================================
\appendix
\section{Side-by-side comparison with the published figures and tables}
\label{app:compare}
% =============================================================================

For visual verification, each panel below pairs an artifact reproduced in this
report \textbf{(a, left)} with the corresponding artifact from the published
paper \textbf{(b, right)}. Published panels are reproduced from
\citet{sheu2024carousel} for comparison. Headline agreement is restated in each
caption; the full discussion is in \S\ref{sec:forward}--\S\ref{sec:recovery}.

\begin{figure}[ht]
\centering
\begin{minipage}[b]{0.45\linewidth}\centering
  \includegraphics[width=\linewidth]{figures/ours_sheu2024b_magnification.png}\\[3pt]
  {\small\textbf{(a) This work}}
\end{minipage}\hfill
\begin{minipage}[b]{0.42\linewidth}\centering
  \includegraphics[width=\linewidth]{figures/orig_sheu2024b_fig4_magnification.png}\\[3pt]
  {\small\textbf{(b) \citet{sheu2024carousel}, Fig.~4}}
\end{minipage}
\caption{\textbf{Magnification map.} (a)~This work: the magnification map at
$\zs=1.432$ of the forward-validated model (the right panel of report
Figure~\ref{fig:forward}; green crosses are the multiple images of a test
source). (b)~The published $\det(A^{-1})$ magnification panel (bottom-middle of
their Fig.~4). The primary tangential critical-curve radius agrees to
$0.5\,\%$ ($\thetaE = 12.96''$ vs $13.03''$). Panel (b) reproduced from
\citet{sheu2024carousel} for comparison.}
\label{fig:cmp-sheu2024b-magnification}
\end{figure}

\begin{figure}[ht]
\centering
\begin{minipage}[b]{0.49\linewidth}\centering
  \includegraphics[width=\linewidth]{figures/predicted_images.png}\\[3pt]
  {\small\textbf{(a) This work}}
\end{minipage}\hfill
\begin{minipage}[b]{0.49\linewidth}\centering
  \includegraphics[width=\linewidth]{figures/orig_sheu2024b_fig6_predicted.png}\\[3pt]
  {\small\textbf{(b) \citet{sheu2024carousel}, Fig.~6}}
\end{minipage}
\caption{\textbf{Predicted images and critical curves.} (a)~This work:
published best-fit model predicted image configurations for four source planes
on the HST F140W image, with the $13.03''$ tangential ring (green dashed)
tracing the observed giant arc. (b)~The published predicted image locations and
critical curves for the same system. Panel (b) reproduced from
\citet{sheu2024carousel}, their Fig.~6, for comparison.}
\label{fig:cmp-sheu2024b-predicted}
\end{figure}

\begin{figure}[ht]
\centering
\begin{minipage}[b]{0.49\linewidth}\centering
  \includegraphics[width=\linewidth]{figures/corner.png}\\[3pt]
  {\small\textbf{(a) This work}}
\end{minipage}\hfill
\begin{minipage}[b]{0.46\linewidth}\centering
  \includegraphics[width=\linewidth]{figures/orig_sheu2024b_fig5_corner.png}\\[3pt]
  {\small\textbf{(b) \citet{sheu2024carousel}, Fig.~5}}
\end{minipage}
\caption{\textbf{Parameter posteriors.} (a)~This work: \texttt{emcee} posterior
for the six free mass-model parameters; red lines mark the published best-fit.
(b)~The published posteriors of the lens-model parameters. $\thetaE^a$ and the
slope $\gammaslope^a$ sit close to the published values. Panel (b) reproduced
from \citet{sheu2024carousel}, their Fig.~5, for comparison.}
\label{fig:cmp-sheu2024b-corner}
\end{figure}

\begin{figure}[ht]
\centering
{\small\textbf{(a) This work (Table~\ref{tab:forward})}}\\[3pt]
\includegraphics[width=0.96\linewidth]{figures/ours_sheu2024b_tab_forward.png}\\[12pt]
{\small\textbf{(b) \citet{sheu2024carousel}, Table~2}}\\[3pt]
\includegraphics[width=0.96\linewidth]{figures/orig_sheu2024b_table2.png}
\caption{\textbf{Forward-model validation table.} (a)~This work: derived
quantities from feeding the published parameters into the multi-plane
\texttt{LensModel}, against the published values. (b)~The published lens-model
parameter table. $\thetaE$ and $M(<\thetaE)$ agree to $0.5\,\%$ and $3.4\,\%$.
Panel (b) reproduced from \citet{sheu2024carousel}, their Table~2, for
comparison.}
\label{fig:cmp-sheu2024b-tabforward}
\end{figure}

\begin{figure}[ht]
\centering
{\small\textbf{(a) This work (Table~\ref{tab:recovery})}}\\[3pt]
\includegraphics[width=0.80\linewidth]{figures/ours_sheu2024b_tab_recovery.png}\\[12pt]
{\small\textbf{(b) \citet{sheu2024carousel}, Table~2}}\\[3pt]
\includegraphics[width=0.96\linewidth]{figures/orig_sheu2024b_table2.png}
\caption{\textbf{Independent-recovery table.} (a)~This work: the PSO\,+\,\texttt{emcee}
recovered mass-model parameters (closed-loop on conjugate image families)
against the published best-fit. (b)~The published lens-model parameter table
(same as above). The headline $\thetaE$ and slope $\gammaslope$ recover in the
right regime ($13.20''$ vs $13.03''$; $1.53$ vs $1.67$); axis ratio, PA, and
shear show the expected ellipticity--shear degeneracy. Panel (b) reproduced
from \citet{sheu2024carousel}, their Table~2, for comparison.}
\label{fig:cmp-sheu2024b-tabrecovery}
\end{figure}

\bibliographystyle{plainnat}
\bibliography{references}

\end{document}
